\chapter{Materials and Methods}

\section{DG-CA3 neural network model}

Based primarily on the models of~\citeonline{kopsickFormation2024, kimAdult2024, yangDynamic2025, chavlisDendrites2017}, the
DG-CA3 network architecture was modeled as shown in Figures~\ref{fig:arquitetura-rede} and~\ref{fig:graph} at a scale of
$\frac{1}{500}$ of the rat hippocampus, with the number of neurons per population described in
Table~\ref{tab:neuron_counts}. Neurons were modeled simply, with a single compartment, but with high biological plausibility,
using neuron and synapse models and cell count and connectivity data from Hippocampome.org, an open-access knowledge base that
compiles multiple data sources on the hippocampal formation~\cite{wheelerHippocampomeorg2023}.

\begin{figure}
    \centering
    \caption{DG-CA3 network architecture. Inhibitory synapses are represented by circles and excitatory ones by arrows.
        The network input is composed of cells from the entorhinal cortex (EC). The dentate gyrus (DG) has the perforant path
        (PP) and the dendrites of its cells in the molecular layer; the granular layer is composed of immature and mature
        Granule Cells (iGCs and mGCs, respectively) and Basket Cells (BCs); the hilus is composed of Mossy Cells (MCs) and
        Hilar Perforant Path-associated cells (HIPP). CA3 is composed of pyramidal CA3 cells (pCA3) and inhibitory CA3 cells
        (iCA3).}
    \includegraphics[width=\textwidth]{figures/arquitetura-rede.png}
    \label{fig:arquitetura-rede}
\end{figure}

\begin{figure}
    \centering
    \caption{Ring graph of the DG-CA3 network. Each node represents a neuron in the model, with the same color scheme as
    Figure~\ref{fig:arquitetura-rede} for each neuron. Each edge represents a synapse, colored according to the presynaptic
    neuron. The EC is represented by the outermost black circle; the lamellar organization of the DG can be seen with mGCs and
    iGCs (red and pink, respectively) and pCA3 cells at the center in teal. In total, the graph contains 3,260 nodes and
    266,596 edges.}
    \includegraphics[width=\textwidth]{figures/network_color_edges}
    \label{fig:graph}
\end{figure}

The network input is composed of cells from the entorhinal cortex (EC), with a total of $N_{EC} = 400$
neurons~\cite{amaralChapter1990, kimAdult2024}. In each simulation, the EC as a whole presents a specific pattern, where each
pattern is represented by an active subpopulation of 10\% of EC neurons~\cite{mcnaughtonDead1991}. Inactive neurons not
belonging to the pattern do not spike during the simulation, whereas active neurons spike according to a Poisson distribution
with a spike rate of $\lambda = \SI{40}{\hertz}$. EC neurons project their axons via the perforant path to neurons with
dendrites in the molecular layer of the DG: Granule Cells (GCs), Mossy Cells (MCs)~\cite{scharfmanHilar2013}, and Basket
Cells (BCs); as well as to CA3 neurons.

GCs are subdivided into two subpopulations: mature Granule Cells (mGCs) and immature Granule Cells (iGCs), representing GCs
generated by adult neurogenesis at their characteristic electrophysiological phase of 4--6 weeks of age~\cite{aimoneRegulation2014}.
In total, the network is composed of $N_{GC} = 2000$ GCs ($\frac{1}{500}$ of the $10^6$ rat granule cells)~\cite{westUnbiased1991},
with 5\% of them being iGCs~\cite{cameronAdult2001}, giving $N_{mGC} = 1900$ and $N_{iGC} = 100$. Following the lamellar
organization of the DG~\cite{sloviterUpdating2012}, GCs are distributed across 20 lamellae, with 100 cells per lamella. Each GC
connects with the BCs, MCs, and CA3 neurons of the same lamella, and also forms random, non-lamellar excitatory connections with
Hilar Perforant Path-associated cells (HIPP).

iGCs, unlike mGCs, possess distinct electrophysiological characteristics that render them more excitable than mGCs. Furthermore,
to simulate their not-yet-complete circuit integration, iGCs receive EC afferent inputs at only a fraction $x_{EC:iGC}$ of the
connection probability between the EC and mGCs. Hippocampome.org~\cite{wheelerHippocampomeorg2023} does not provide iGC
characteristics. Therefore, the parameters used for iGCs were fitted from the spiking data
of~\citeonline{espositoNeuronal2005} using the method of~\citeonline{nelderSimplex1965}
(Table~\ref{tab:izhikevich_neuron_params}).

BCs ($N_{BC} = 40$) are inhibitory cells that ensure sparse GC activation through winner-take-all competition among
them~\cite{coultripCortical1992, chavlisDendrites2017, kimAdult2024}, where only the most active GCs in a lamella remain active,
inhibiting the rest. HIPPs ($N_{HIPP} = 60$) also contribute to the sparsity of GC activation, with a more global inhibition
acting on all GCs in a lamella. MCs ($N_{MC} = 100$) are excitatory hilar cells that receive excitation from GCs and project
back to GCs, BCs, and HIPPs via interlamellar connections. Although these excitatory MC-to-GC projections exist, their net
effect is generally inhibitory through the control of BCs and HIPPs~\cite{myersRole2009, scharfmanHilar2013}.

The HIPP parameters in Hippocampome.org~\cite{wheelerHippocampomeorg2023} model the ``rebound burst spiking'' characteristic,
in which the neuron produces a burst of spikes following the end of a hyperpolarization period. This model is incompatible with
the network, as large hyperpolarizations lead to infinite spike bursts. Therefore, the parameters used for HIPPs were taken
from~\citeonline{modakIzhikevich2018}.

CA3, unlike the DG, does not follow a lamellar structure in the model, as it forms a far more integrative neural
network~\cite{pakHippocampal2022, watsonHuman2025}. The CA3 model comprises 600 pyramidal cells ($N_{pCA3} = 600$) and 60
inhibitory cells ($N_{iCA3} = 60$), modeled according to the physiological data for CA3 basket cells~\cite{wheelerHippocampomeorg2023}
so as to simplify the variety of inhibitory neurons present in CA3~\cite{kopsickFormation2024}. Both neuronal populations
receive afferent inputs from the EC and from GCs, with iCA3 cells inhibiting pCA3 cells. pCA3 cells excite iCA3 cells and,
importantly, form a recurrent network by connecting with other pCA3 cells, a fundamental characteristic for the autoassociation
and pattern completion functions of CA3~\cite{kopsickFormation2024, rollsMechanisms2013}. pCA3 cells also send feedback
projections to the DG, connecting with MCs in the model, a process that contributes to pattern
separation~\cite{myersPattern2011}.

All simulations were performed with Brian2~\cite{stimbergBrian2019a}, using the 4th-order Runge-Kutta method with a fixed time
step of \SI{0.1}{\milli\second}~\cite{butcherHistory1996}. Neuron and synapse parameters can be found in
Tables~\ref{tab:izhikevich_neuron_params} and~\ref{tab:synapse_params}, respectively. The duration of a single simulation for
one input pattern is \SI{500}{\milli\second} for the network to reach a steady state, with the pattern recording phase
occurring subsequently for \SI{1000}{\milli\second}. Synaptic connections are kept fixed across all simulations.

\input{tabelas/neuron_counts.tex}

\section{Neuron model}

Neurons were modeled according to the 9-parameter Izhikevich neuron model~\cite[cap.~8]{izhikevichDynamical2006} with a single
compartment, without considering dendrites or axons. This model was chosen in this work and by~\citeonline{wheelerHippocampomeorg2023}
for being capable of capturing the dynamic behavior of neurons across a wide variety of conditions with biological plausibility,
comparable to the Hodgkin-Huxley model~\cite{hodgkinQuantitative1952b}, while presenting a simpler and computationally more
efficient mathematical formulation. The spike traces of neurons can be seen in Figure~\ref{fig:spike_traces}. The Izhikevich
neuron model is described by the following equations:


\begin{equation}
\label{eq_izhikevich_1}
C_m \frac{dV_m}{dt} = k (V_m - V_r)(V_m - V_t) - u + I
\end{equation}

\begin{equation}
\label{eq_izhikevich_2}
\frac{du}{dt} = a [b(V_m-V_r) - u]
\end{equation}

Where $V_m$ is the membrane potential, $u$ is the recovery variable, $C_m$ is the membrane capacitance, $V_r$ is the resting
potential, $V_t$ is the threshold potential, $I$ is the total current flowing into the neuron, and $k$, $a$, and $b$ are
constants that define the dynamic characteristics of the neuron. In addition to the differential equations above, which define
the temporal evolution of the membrane potential and recovery variable, the Izhikevich neuron model also includes a rule for
action potential generation, defined by Equation~\ref{eq_izhikevich_3}.

\begin{equation}
\label{eq_izhikevich_3}
\text{if } V_m \geq V_{\text{peak}}, \quad
\begin{cases}
V_m \gets V_{min} \\
u \gets u + d
\end{cases}
\end{equation}

When the membrane potential reaches the peak value $V_{\text{peak}}$, an action potential is generated, the membrane potential
is reset to the after-spike potential $V_{min}$, and the recovery variable $u$ is incremented by $d$, making subsequent spike
generation more difficult.


\begin{figure}[H]
    \centering
    \caption{Spike traces of neurons under constant current as indicated by each subfigure title.}
    \includegraphics[width=0.9\textwidth]{figures/spike_traces}
    \label{fig:spike_traces}
\end{figure}

\input{tabelas/izhikevich.tex}

\section{Synapse model}

The synapse model, like the neuron model, was defined based on Hippocampome.org~\cite{wheelerHippocampomeorg2023}, following the
formulation of~\citeonline{sennAlgorithm2001}\cite{mongilloSynaptic2008a}. This model captures short-term plasticity, whether
short-term depression, caused by neurotransmitter depletion, or short-term potentiation, caused by calcium
accumulation, both occurring on the scale of tenths of seconds. Each synapse has 5 parameters (described in
Table~\ref{tab:synapse_params}): the maximum synaptic conductance in the absence of any synaptic resource depletion $g$, the
proportion of resources utilized per spike $U_{se}$, the synaptic current decay time constant $\tau_d$, the facilitation time
constant $\tau_f$, and the resource recovery time constant $\tau_r$~\cite{moradiNormalized2022}. Another parameter shown in
Table~\ref{tab:synapse_params} is the connection probability $P$. Although Hippocampome.org provides data on connection
probabilities between neuronal populations \textit{in vivo}, the reduced scale of the network necessitated increasing these
probabilities so that network activity could be maintained while sparse activity was preserved.

The model is described by three state variables: synaptic resource utilization ($U$), resource recovery ($R$), and the
percentage of resources in an active state ($A$). Initially, $U_{t_0} = 0$, $R_{t_0} = 1$, and $A_{t_0} = 0$, since all
resources are available for use. The temporal evolution of these variables is governed by the following system of differential
equations:

\begin{equation}
    \label{eq_tsodyks_dU}
    \frac{dU}{dt} = \frac{-U}{\tau_f} + U_{se}(1-U_{-}) \delta(\Delta t_i)
\end{equation}

\begin{equation}
    \label{eq_tsodyks_dR}
    \frac{dR}{dt} = \frac{1-R-A}{\tau_r} - U_{+} R_{-} \delta(\Delta t_i)
\end{equation}

\begin{equation}
    \label{eq_tsodyks_dA}
    \frac{dA}{dt} = \frac{-A}{\tau_d} + U_{+} R_{-} \delta(\Delta t_i)
\end{equation}

where $\delta$ is the Dirac delta function, yielding 1 only when $\Delta t_i = t - t_i = 0$, i.e., only at time $t$
corresponding to the synaptic event time $t_i$. $U_{+}$ corresponds to the value of $U$ immediately after the synaptic event,
while $R_{-}$ corresponds to the value of $R$ immediately before the same event.

From these equations, the synaptic current is given by:

\begin{equation}
    \label{eq_tsodyks_I}
    I = k \cdot A \cdot g \cdot (V_m - E)
\end{equation}

where $V_m$ is the postsynaptic membrane potential, $E$ is the synaptic reversal potential, for inhibitory and excitatory
synapses, respectively, $E_{inh} = \SI{-86}{\milli\volt}$ and $E_{exc} = \SI{0}{\milli\volt}$, and $k$ is a scaling constant
defined as $k = 10$ for all synapses. This scaling constant is necessary due to the reduced scale of the network, since, given
the low number of synapses in the model compared to the rat hippocampus, the entire network would fall silent without this
scaling. All synaptic currents $I$ of a neuron are summed and affect the membrane potential $V_m$ according to
Equation~\ref{eq_izhikevich_1}.


\input{tabelas/sinapses.tex}


\section{Long-term plasticity}

Long-term plasticity will be implemented exclusively in the recurrent synapses between pyramidal CA3 cells (pCA3), given their
fundamental role in the formation of neural assemblies and to ensure the computational efficiency of the model. The plasticity
model will be based on the work of~\citeonline{kopsickFormation2024}, which employs a Spike-Timing-Dependent Plasticity (STDP)
rule.

Specifically, a symmetric, Hebbian STDP rule will be used, which promotes the strengthening of synapses between neurons that
spike in close temporal proximity, regardless of spiking order. This rule is described by the following equation:

\begin{equation}
    \label{eq:stdp}
    \Delta w = A e^{-|\Delta t|/\tau}
\end{equation}

Where $\Delta w$ is the change in synaptic weight, $A$ is a parameter that determines the maximum magnitude of the weight
change, $\tau$ is the plasticity decay time constant, and $\Delta t$ is the temporal difference between the spikes of the pre-
and postsynaptic neuron. Following the model of~\cite{kopsickFormation2024}, the time constant $\tau$ will be set to
\SI{20}{\milli\second}. The parameter $A$ will be adjusted to modulate the learning rate of the network.

\begin{figure}[H]
    \centering
    \caption{Symmetric Hebbian STDP rule with $A = 70$.}
    \includegraphics[width=0.5\textwidth]{figures/symmetric_stdp}
    \label{fig:symmetric_stdp}
\end{figure}

To prevent synaptic saturation, a synaptic weight renormalization rule (Section~\ref{sec:protocolo_treinamento_teste}) and a
maximum synaptic weight $w_{max}$ to be determined experimentally will be employed.

\section{Temporal neurogenesis}

To investigate the functional impact of the continuous integration of new neurons, a temporal neurogenesis model will be
implemented. This process is fundamental for the maintenance of hippocampal functions such as learning and memory over
time~\cite{aimoneRegulation2014, berdugo-vegaSharpening2023}. The simulation will be divided into two main phases. In the first
phase, the network will operate with its initial configuration, with $N_{mGC} = 1900$ and $N_{iGC} = 100$. During this phase,
the network will be exposed to a set of $N$ input patterns, which will be learned and stored in the recurrent synapses of CA3
via the STDP mechanism. The number of input patterns $N$ will be varied experimentally to determine how many patterns the model
can store without a significant deficit in pattern completion (see Section~\ref{sec:quantificacao}).

After the initial learning phase, iGCs will undergo a discrete simulated maturation process, transforming into mGCs all at
once. This transition reflects the biological changes that occur as new neurons become fully integrated into the DG circuit,
transitioning from a hyperexcitable state to a more stable mature state~\cite{abbottAdult2020}. Maturation will be implemented
by updating their electrophysiological and synaptic parameters to the values corresponding to those of mGCs
(Tables~\ref{tab:izhikevich_neuron_params} and~\ref{tab:synapse_params}). Existing efferent connections, formed during the
immature phase, will be preserved. Additionally, the newly matured cells will have their EC afferent connectivity increased to
the same level as mGCs, completing their functional integration into the circuit.

In the second phase of the simulation, following the maturation of the first iGCs, 100 new iGCs will be added, simulating the
continuous integration of new neurons by adult neurogenesis, and another batch of $N$ patterns will be presented to the
network. This process will be repeated 3 times, making it possible to evaluate network performance progressively, tracking the
maturation of iGCs. Two control models will be implemented: one in which only mGCs are added to the network without iGCs, and
another in which, when iGCs are added, an equivalent number of mGCs undergo neuronal death. The purpose of these control models
is to isolate the role of adult neurogenesis and of the immature cells it generates, comparing against the effects of control
models that also feature a gradual increase in neuron count.


\section{Pattern separation}\label{sec:separacao_padroes}

The methodology for quantifying pattern separation was based on that used by~\citeonline{kimEffect2024}. To analyze pattern
separation under different levels of input pattern similarity and the impact of adult neurogenesis, 11 different model variants
were tested: one control without adult neurogenesis ($N_{mGC} = 2000$) and 10 others with neurogenesis at different
connectivity levels ($N_{mGC} = 1900$ and $N_{iGC} = 100$), as illustrated in Figure~\ref{fig:diagrama_entrada_saida}. The
connectivity level is expressed as a percentage, where 100\% corresponds to iGCs having the same EC-GC connection probability
as mGCs; the intent is to simulate the incomplete circuit integration of iGCs, whose dendrites are not yet fully developed. For
each model, 20 sets of 10 input patterns are generated. Each set consists of one randomly generated original pattern and 9
patterns similar to it, ranging from 90\% to 10\% similarity. An input pattern consists of a binary vector of size
$N_{EC} = 400$, representing the activation state of each EC neuron, in which active neurons spike according to a Poisson
distribution at a frequency of \SI{40}{\hertz}. Every input pattern has exactly 10\% active neurons. To generate patterns with
different levels of similarity, a random fraction of $X\%$ of the active neurons is kept active, while the remaining
$(100 - X)\%$ are deactivated and an equal number of initially inactive neurons is activated, to maintain the 10\% activity
level (40 active neurons).

Each pattern is simulated for \SI{1500}{\milli\second} of simulation time, with the first \SI{500}{\milli\second} discarded
while the network stabilizes. The output pattern of a neuronal population, whether mGCs, iGCs, or pCA3 cells, consists of a
binary vector with value 1 for the index of each neuron that spiked at least once during the \SI{1000}{\milli\second} interval.

\begin{figure}[H]
    \centering
    \caption{Diagram of the input and output pattern generation and recall protocol.
    \textbf{A.} Input patterns are generated in sets of 10 patterns of size $N_{EC} = 400$; the first is generated randomly
    and the remaining ones are generated from it at different levels of similarity by randomly activating and deactivating
    neurons.
    \textbf{B.} 11 model variants are simulated: one without neurogenesis and 10 with neurogenesis at different connectivity
    levels, from 100\% down to 10\% connectivity.
    \textbf{C.} Spike trace of EC neurons during a simulation; each generated pattern corresponds to the activation of EC
    neurons, with active neurons spiking at a frequency of \SI{40}{\hertz}; each of the 11 models is exposed to 20 sets of
    10 input patterns, totaling 200 simulations of \SI{1500}{\milli\second} each.
    \textbf{D.} Spike trace of mGCs during a simulation; the first \SI{500}{\milli\second} are discarded while the network
    stabilizes; the remainder corresponds to the output pattern.
    \textbf{E.} The output pattern, illustrated here for the mGC population, consists of a binary vector with value 1 for the
    index of each neuron that spiked at least once during the \SI{1000}{\milli\second} interval.}
    \includegraphics[width=\textwidth]{figures/diagrama_entrada_saida}
    \label{fig:diagrama_entrada_saida}
\end{figure}

To characterize pattern separation, the overlap between the neural activity patterns at the input (entorhinal cortex cells) and
output (DG Granule Cells or CA3 pyramidal cells) of the network is compared. A pattern is defined by a binary representation of
size $N$, where $N$ is the total number of neurons in a specific population, with neurons that spiked at least once during the
simulation time interval (\SI{1000}{\milli\second}) represented by 1 and those that did not by 0. For a pair of patterns
$A^{(l)}$ and $B^{(l)}$ (where $l \in \{in, out\}$ for input and output, respectively), the distance between patterns
$D_p^{(l)}$ is defined as:

\begin{equation}
    \label{eq:dp}
    D_p^{(l)} = \frac{O^{(l)}}{D_a^{(l)}}
\end{equation}

In this equation, $O^{(l)}$ represents the degree of orthogonalization and $D_a^{(l)}$ the mean degree of activation of the
two patterns. The mean degree of activation $D_a^{(l)}$ is the arithmetic mean of the degrees of activation of each pattern,
$A^{(l)}$ and $B^{(l)}$:

\begin{equation}
    \label{eq:da}
    D_a^{(l)} = \frac{D_a^{(A^{(l)})} + D_a^{(B^{(l)})}}{2}
\end{equation}

The degree of activation of an individual pattern is the fraction of active neurons (represented by 1 in a binary encoding) in
the pattern. The orthogonalization degree $O^{(l)}$, which measures the dissimilarity between patterns, is calculated from the
Pearson correlation coefficient, $\rho^{(l)}$:

\begin{equation}
    \label{eq:o}
    O^{(l)} = \frac{1 - \rho^{(l)}}{2}
\end{equation}

Where $\rho^{(l)}$ is the Pearson correlation coefficient between patterns $A^{(l)}$ and $B^{(l)}$.
Considering $\{a_i^{(l)}\}$ and $\{b_i^{(l)}\}$ ($i=1, \dots, N_l$) as the binary representations of the state of the
$i$-th cell in patterns $A^{(l)}$ and $B^{(l)}$ ($l \in \{in, out\}$), the Pearson correlation coefficient is given by:

\begin{equation}
    \label{eq:pearson}
    \rho^{(l)} = \frac{\sum_{i=1}^{N_l} \Delta a_i^{(l)} \cdot \Delta b_i^{(l)}}{\sqrt{\sum_{i=1}^{N_l} (\Delta a_i^{(l)})^2} \sqrt{\sum_{i=1}^{N_l} (\Delta b_i^{(l)})^2}}
\end{equation}

where $\Delta a_i^{(l)} = a_i^{(l)} - \langle a^{(l)} \rangle$ and $\Delta b_i^{(l)} = b_i^{(l)} - \langle b^{(l)} \rangle$.
The notation $\langle \dots \rangle$ indicates the population mean over all cells. The value of $\rho^{(l)}$ varies between
-1 and 1, therefore $O^{(l)}$ varies between 0 and 1 (Equation~\ref{eq:o}).

From the input ($D_p^{(in)}$) and output ($D_p^{(out)}$) pattern distances, the degree of pattern separation,
$\mathcal{S}_D$, is calculated as the ratio between them:

\begin{equation}
    \label{eq:sd}
    \mathcal{S}_D = \frac{D_p^{(out)}}{D_p^{(in)}}
\end{equation}

A value of $\mathcal{S}_D > 1$ indicates that the output patterns are more distinct than the input patterns, characterizing
pattern separation. Conversely, $\mathcal{S}_D < 1$ indicates pattern convergence, where the output patterns become more
similar to each other.

\section{Autoassociation and pattern completion}

\subsection{Training and testing protocol}\label{sec:protocolo_treinamento_teste}

The network training and testing protocol for evaluating pattern storage and retrieval will be based on the method described
by~\citeonline{kopsickFormation2024}, who implemented a full-scale CA3 model also using Hippocampome.org data and developed a
biologically plausible training protocol.

During the training phase, the network will be exposed to a set of distinct input patterns. Each pattern consists of the
activation of a specific subpopulation of pCA3 cells via current injection. This activation induces a spike train in a
\SI{20}{\milli\second} time window, corresponding to one gamma cycle. The different patterns will be presented in sequence,
separated by \SI{200}{\milli\second} windows, simulating a theta-gamma neural code. Throughout the training phase, synaptic
plasticity (STDP) will be enabled, allowing the strengthening of connections between neurons encoding the same pattern and thus
forming a neural assembly. After a number of pattern presentations to be determined experimentally, the synaptic weights between
pyramidal cells will be renormalized, a process that simulates the synaptic homeostasis occurring during slow-wave sleep, to
prevent synaptic saturation and stabilize the learned patterns~\cite{gonzalez-ruedaActivityDependent2018, kopsickFormation2024}.
This renormalization consists of dividing each synaptic weight between pyramidal cells by the same factor, such that the mean
weight across all synapses returns to its initial value. This procedure preserves the relative weight distribution, thereby
maintaining the learned associations without allowing synapses to saturate.

During the testing phase, synaptic plasticity will be disabled in order to assess the network's ability to retrieve the stored
patterns without modification to the connections. Degraded versions of the original patterns will be presented, in which only a
fraction of the neurons in the corresponding assembly is activated. The network response will then be analyzed as described in
Section~\ref{sec:quantificacao} to verify whether the activity of the recurrent connections can reconstruct the complete
pattern through pattern completion.

\subsection{Quantification}\label{sec:quantificacao}

The strength of autoassociation, that is, the formation of assemblies, will be measured through two characteristics of the
synapses between pyramidal cells. The first is the signal-to-noise ratio (SNR) of autoassociation, defined as the mean synaptic
weight between neurons of the same assembly divided by the mean synaptic weight between neurons of different assemblies. A high
SNR indicates a strong distinction between intra- and inter-assembly connections. The second metric is the percentage of
synapses within an assembly that reached the maximum weight, which serves as an indicator of synaptic saturation.

Pattern completion will be assessed through a pattern reconstruction accuracy metric, based
on~\cite{kopsickFormation2024}. This metric uses the Pearson correlation coefficient (Equation~\ref{eq:pearson}) between input
or output patterns during the training and testing phases.

First, the input pattern correlation, $\rho^{(in)}$, is calculated, which is the correlation between the complete EC pattern
presented during the training phase and the degraded, or incomplete, pattern presented during the testing phase. Next, the
output pattern correlation, $\rho^{(out)}$, is calculated, corresponding to the correlation between the pCA3 neural activity
in response to the complete training pattern during assembly encoding and the activity in response to the degraded test pattern.

The pattern reconstruction accuracy, $R_p$, is then calculated as:

\begin{equation}
    \label{eq:rp}
    R_p = \frac{\rho^{(out)} - \rho^{(in)}}{1 - \rho^{(in)}}
\end{equation}

A value of $R_p > 0$ indicates that pattern completion occurred, i.e., the output representation became more similar to the
original pattern than the degraded input. A value of $R_p = 1$ represents perfect reconstruction.
